% =============================================================================
% Vertical A0 Scientific Poster
% Title: 3D Organ Segmentation from CT Scans
% Framework: beamerposter + tcolorbox | Size: A0 (841 x 1189 mm) | 3 columns
% UTAD Brand Colors: navy RGB(25,34,71) | red RGB(190,29,44) | gray RGB(198,198,198)
% Compile with: pdflatex poster.tex
% =============================================================================
\documentclass[final,t]{beamer}

% A0 poster — portrait (841 × 1189 mm). scale=1.4 → ~15 pt body.
% orientation=portrait must be explicit; default is landscape.
% Do NOT override \textheight: beamer measures the headline/footline
% templates and computes frametextheight = textheight - head - foot.
% Overriding causes double-subtraction and a 14 cm frame-too-small bug.
\usepackage[size=a0,orientation=portrait,scale=1.4]{beamerposter}

% --------------------------------------------------------------------------
% Packages
% --------------------------------------------------------------------------
\usepackage{lmodern}
\usepackage{amsmath,amssymb}
\usepackage{graphicx}
\usepackage{booktabs}
\usepackage{colortbl}          % \rowcolor in tables (avoids xcolor option conflict)
\usepackage{array}
\usepackage{fontawesome5}
\usepackage{tikz}
\usetikzlibrary{shapes.geometric,arrows.meta,positioning,fit,backgrounds,calc}
\usepackage{tcolorbox}
\tcbuselibrary{skins,breakable}
\usepackage{qrcode}
\usepackage{ragged2e}          % \RaggedRight, \justifying
\usepackage{calc}
\usepackage{url}

% --------------------------------------------------------------------------
% UTAD Brand Colors
% --------------------------------------------------------------------------
\definecolor{utadnavy}{RGB}{25,34,71}
\definecolor{utadred}{RGB}{190,29,44}
\definecolor{utadgray}{RGB}{198,198,198}
\definecolor{utadlightbg}{RGB}{237,241,250}

% Semantic colors for architecture comparison
\definecolor{cnnblue}{RGB}{41,128,185}
\definecolor{transgreen}{RGB}{39,174,96}
\definecolor{foundorange}{RGB}{200,90,0}

% --------------------------------------------------------------------------
% Beamer Setup (minimal)
% --------------------------------------------------------------------------
\usetheme{default}
\setbeamertemplate{navigation symbols}{}
\setbeamertemplate{frametitle}{}
\setbeamercolor{background canvas}{bg=white}

% --------------------------------------------------------------------------
% tcolorbox: poster section block
% --------------------------------------------------------------------------
\newtcolorbox{pblock}[1]{%
  title={#1},
  colbacktitle=utadnavy,
  coltitle=white,
  colback=utadlightbg,
  colframe=utadnavy,
  fonttitle=\large\bfseries,
  boxrule=1.5pt,
  arc=5pt,
  lefttitle=7pt,
  toptitle=2pt,
  bottomtitle=2pt,
  top=4pt, bottom=4pt, left=9pt, right=9pt,
  before={\vskip2pt},
  after={\vskip2pt},
  enhanced,
  drop shadow={utadnavy!25!white,opacity=0.5},
}

% --------------------------------------------------------------------------
% tcolorbox: key-stat chip (used in 2×2 grid)
% --------------------------------------------------------------------------
\newtcolorbox{statchip}[1]{%
  title={#1},
  colbacktitle=utadnavy,
  coltitle=white,
  colback=utadred!88!black,
  colframe=utadnavy,
  colupper=white,
  fonttitle=\small\bfseries,
  fontupper=\huge\bfseries,
  boxrule=0pt,
  arc=8pt,
  lefttitle=0pt,
  toptitle=3pt,
  bottomtitle=3pt,
  halign=center,
  halign title=center,
  height=5.5cm,
  valign=center,
  top=0pt, bottom=0pt, left=4pt, right=4pt,
  before={\vskip0pt},
  after={\vskip0pt},
  enhanced,
}

% --------------------------------------------------------------------------
% Highlight box (inline note)
% --------------------------------------------------------------------------
\newtcolorbox{notebox}[1][utadnavy]{%
  colback=#1!8,
  colframe=#1,
  arc=4pt,
  boxrule=1pt,
  top=4pt, bottom=4pt, left=7pt, right=7pt,
  before={\vskip4pt},
  after={\vskip4pt},
}

% --------------------------------------------------------------------------
% Headline (title banner)  — A0 = 84.1 cm wide
% Layout:  [logos 20cm] [2cm gap] [title 40cm centered] [2cm gap] [QR 18cm]
% --------------------------------------------------------------------------
\setbeamertemplate{headline}{%
  \leavevmode
  \begin{beamercolorbox}[wd=\paperwidth]{headline}
  \begin{tikzpicture}
    % Background fill — 18 cm tall header
    \fill[utadnavy] (0,0) rectangle (\paperwidth,7cm);
    % Thin red accent at the bottom of the header
    \fill[utadred]  (0,0) rectangle (\paperwidth,0.6cm);

    % ---- Left: Institution logos ----------------------------------------
    % TODO: Replace fcolorbox placeholders with actual logo files.
    %       Create a logos/ subfolder next to this file and add:
    %         logos/champlain.png  logos/utad.png  logos/isec.png  logos/docdo.png
    %       Then replace each \fcolorbox block with \includegraphics[height=3cm]{logos/xx.png}
    % \node[anchor=north west, inner sep=0pt] at (1.5cm, 17.5cm) {%
    %   \begin{minipage}[c][16cm][c]{19cm}
    %     \centering
    %     \fcolorbox{utadgray}{utadnavy!70!black}{%
    %       \parbox[c][3cm][c]{4.2cm}{\centering\color{white}\footnotesize\bfseries
    %         Champlain College}}\\[4pt]
    %     \fcolorbox{utadgray}{utadnavy!70!black}{%
    %       \parbox[c][3cm][c]{4.2cm}{\centering\color{white}\footnotesize\bfseries
    %         UTAD\\Portugal}}\\[4pt]
    %     \fcolorbox{utadgray}{utadnavy!70!black}{%
    %       \parbox[c][3cm][c]{4.2cm}{\centering\color{white}\footnotesize\bfseries
    %         ISEC\\Coimbra}}\\[4pt]
    %     \fcolorbox{utadgray}{utadnavy!70!black}{%
    %       \parbox[c][3cm][c]{4.2cm}{\centering\color{white}\footnotesize\bfseries
    %         DocDo Corp.}}
    %   \end{minipage}%
    % };

    % ---- Center: Title, subtitle, authors, affiliations ------------------
    \node[anchor=north, inner sep=0pt, text width=44cm, align=center]
      at (0.5\paperwidth, 6cm) {%
        {\color{white}\fontsize{52}{60}\selectfont\bfseries
          3D Organ Segmentation from CT Scans}\\[6pt]
        {\color{utadgray}\fontsize{29}{36}\selectfont
          An Agentic Structured Survey of Deep Learning
          Approaches for Surgical Planning}\\[10pt]
        {\color{white}\fontsize{22}{29}\selectfont
          A.\,T.\,Nogueira$^{1,2,3,4}$\enspace
          M.\,H.\,R.\,Martins$^{4}$\enspace
          N.\,M.\,Ferreira$^{3}$\enspace
          V.\,M.\,Filipe$^{2}$}\\[4pt]
        {\color{utadgray}\fontsize{17}{22}\selectfont
          $^{1}$Champlain College, USA\quad
          $^{2}$UTAD, Portugal\quad
          $^{3}$ISEC, Portugal\quad
          $^{4}$DocDo Corp., USA}
    };

    % ---- Right: QR code -------------------------------------------------
    % Header is 7 cm tall. anchor=center keeps QR fully inside the navy band.
    \node[anchor=center, inner sep=0pt]
      at (\dimexpr\paperwidth-5cm, 3.8cm) {%
        \begin{minipage}[c][4.5cm][c]{4.5cm}
          \centering
          \color{white}
          \qrcode[hyperlink,height=4.5cm]{https://github.com/gameguild-gg/docdo-paper}
        \end{minipage}%
    };

  \end{tikzpicture}%
  \end{beamercolorbox}%
}

% --------------------------------------------------------------------------
% Footline
% --------------------------------------------------------------------------
\setbeamertemplate{footline}{%
  \leavevmode
  \begin{beamercolorbox}[wd=\paperwidth]{footline}
  \begin{tikzpicture}
    \fill[utadnavy] (0,0) rectangle (\paperwidth,3.2cm);
    \fill[utadred]  (0,2.9cm) rectangle (\paperwidth,3.2cm);
    \node[anchor=center] at (0.5\paperwidth, 1.6cm) {%
      {\color{white}\fontsize{19}{25}\selectfont
        \faEnvelope\ tolstenko@gmail.com
        \quad\textbf{|}\quad
        \faGithub\ https://github.com/gameguild-gg/docdo-paper
      }
    };
  \end{tikzpicture}
  \end{beamercolorbox}%
}

% =============================================================================
\begin{document}
\begin{frame}[t]
\vspace{0.05cm}
\begin{columns}[t,totalwidth=\textwidth]

% =============================================================================
%  COLUMN 1 — Motivation, Research Questions, Methodology
% =============================================================================
\begin{column}{0.315\textwidth}

%--- Motivation & Clinical Need ----------------------------------------------
\begin{pblock}{Motivation \& Clinical Need}
\RaggedRight
A gap remains between benchmark performance and clinical deployment
of 3D CT organ segmentation for surgical planning:

\begin{itemize}\setlength\itemsep{2pt}
  \item Manual segmentation: \textbf{15--45 min/structure}
  \item EAU recommends 3D reconstruction for complex renal tumours but cites
        lack of \textit{``standardised, rapid segmentation tools''}
  \item 3D preoperative planning \textbf{reduces positive surgical margins by
        23--47\%} in partial nephrectomy
  \item Survey of N=387 surgeons: \textbf{68\%} cite
        \textit{``time-intensive preparation''} as the main barrier
\end{itemize}

\vspace{0.2em}
% Key-stat chips (2 × 2 grid)
\begin{center}
  \begin{minipage}[t]{0.48\linewidth}
    \begin{statchip}{Studies Reviewed (2017--2025)}
      52
    \end{statchip}
  \end{minipage}\hfill
  \begin{minipage}[t]{0.48\linewidth}
    \begin{statchip}{BTCV Dice --- nnU-Net}
      82.3\%
    \end{statchip}
  \end{minipage}\\[0.5em]
  \begin{minipage}[t]{0.48\linewidth}
    \begin{statchip}{CNN Dominance in Literature}
      98.1\%
    \end{statchip}
  \end{minipage}\hfill
  \begin{minipage}[t]{0.48\linewidth}
    \begin{statchip}{Top BTCV Dice --- STU-Net-L}
      87.1\%
    \end{statchip}
  \end{minipage}
\end{center}
\end{pblock}

%--- Research Questions ------------------------------------------------------
\begin{pblock}{Research Questions}
\begin{enumerate}\setlength\itemsep{5pt}
  \item[\textbf{1}] What \textbf{architectural paradigms} dominate 3D CT organ
        segmentation? How do design principles, inductive biases, and
        computational costs differ?
  \item[\textbf{2}] Which \textbf{datasets and metrics} are standard?
        What are their biases and generalisation implications?
  \item[\textbf{3}] What are the \textbf{accuracy vs.\ cost vs.\ annotation}
        trade-offs? How well do methods generalise across domains?
  \item[\textbf{4}] What \textbf{gaps block clinical deployment}: reproducibility,
        domain shift, regulatory hurdles, integration complexity?
\end{enumerate}
\end{pblock}

%--- Methodology -------------------------------------------------------------
\begin{pblock}{Methodology\newline(PRISMA 2020 Adapted)}
\begin{center}
\resizebox{0.87\columnwidth}{!}{%
\begin{tikzpicture}[node distance=0.35cm,
  mbox/.style={rectangle, draw=utadnavy, fill=utadlightbg,
               text width=10.5cm, minimum height=1.2cm, align=center,
               font=\normalsize\bfseries, rounded corners=3pt, line width=1.3pt,
               inner sep=4pt},
  ebox/.style={rectangle, draw=utadred, fill=utadred!6,
               text width=8.5cm, minimum height=0.9cm, align=center,
               font=\small, rounded corners=3pt, line width=0.9pt,
               inner sep=3pt},
  ibox/.style={rectangle, draw=utadgray!70!black, fill=utadgray!25,
               text width=8.5cm, minimum height=0.9cm, align=center,
               font=\small, rounded corners=3pt, line width=0.9pt,
               inner sep=3pt},
  farrow/.style={->, >=stealth, thick, utadnavy},
  earrow/.style={->, >=stealth, thick, utadred},
  iarrow/.style={->, >=stealth, thick, utadgray!70!black}]

  \node[mbox] (db) {Records Identified (n\,=\,2{,}821)\\
                    \small PubMed, arXiv, Semantic Scholar};
  \node[ebox, right=0.8cm of db] (src) {Search executed Nov\,2025\,--\,Jan\,2026\\
                                         DOI + title dedup,\\
                                         citation tracking on 12 seminal papers};
  \node[mbox, below=of db] (es) {Elasticsearch Pre-filter (n\,=\,638)\\
                                  \small Boolean IC queries on abstracts};
  \node[ebox, right=0.8cm of es] (ex0) {Excluded n\,=\,2{,}183\\
                                         No 3D/organ/DL keyword};
  \node[mbox, below=of es] (ai) {AI Screening: 638 $\to$ 161\\
                                  \small GPT-4o-mini, GPT-5-nano, GPT-5.2};
  \node[ebox, right=0.8cm of ai] (ex1) {Excluded n\,=\,477\\
                                         No 3D DL, non-CT};
  \node[mbox, below=of ai] (ft) {Full-text Assessed (n\,=\,63)\\
                                  \small 98 not retrieved (paywalled / no PDF)};
  \node[ebox, right=0.8cm of ft] (ex2) {Excluded n\,=\,11\\
                                         Final full-text screen};
  \node[mbox, below=of ft,
        fill=utadnavy!12, draw=utadnavy, line width=2pt] (inc)
    {\textbf{Included: n\,=\,52}\\
     \small Synthesis (n=52), Per-organ (n=32)};

  \draw[farrow] (db) -- (es);
  \draw[farrow] (es) -- (ai);
  \draw[farrow] (ai) -- (ft);
  \draw[farrow] (ft) -- (inc);
  \draw[earrow] (es.east)  -- (ex0.west);
  \draw[earrow] (ai.east)  -- (ex1.west);
  \draw[earrow] (ft.east)  -- (ex2.west);
\end{tikzpicture}%
}
\end{center}
\vspace{0.3em}
{\small
  \textbf{Literature window:} 2015 -- 2025\quad
  \textbf{Full-text retrieval:} 39.1\%\\
  \textbf{Human audit:} $\kappa = 0.89$ (10\% stratified sample, n=64,
  two domain experts blind to AI decisions)
}
\end{pblock}

\end{column}
\hfill

% =============================================================================
%  COLUMN 2 — Results, Per-Organ, Architecture
% =============================================================================
\begin{column}{0.315\textwidth}

%--- Key Results: BTCV -------------------------------------------------------
\begin{pblock}{Results: BTCV Multi-Organ Benchmark (13 Organs)}
\vspace{0.4em}
{\small
\begin{tabular}{@{}llccc@{}}
\toprule
\textbf{Method} & \textbf{Type} &
  \textbf{Dice\,\%} & \textbf{HD95\,mm} & \textbf{Pre-train} \\
\midrule
\rowcolor{cnnblue!10}
3D U-Net        & CNN     & 74.2 & 18.4 & ---         \\
\rowcolor{cnnblue!10}
V-Net           & CNN     & 75.8 & 16.7 & ---         \\
\rowcolor{cnnblue!10}
Att.\ U-Net     & CNN     & 77.1 & 14.2 & ---         \\
\rowcolor{cnnblue!20}
nnU-Net         & \textbf{CNN}  & \textbf{82.3} & 21.4 & --- \\
\rowcolor{cnnblue!20}
MedNeXt         & \textbf{CNN}  & \textbf{86.2} & 12.1 & --- \\
\rowcolor{cnnblue!20}
STU-Net-L       & \textbf{CNN}  & \textbf{87.1} & \textbf{9.8} & --- \\
\midrule
\rowcolor{transgreen!10}
UNETR           & Trans.  & 78.7 & 29.1 & \checkmark  \\
\rowcolor{transgreen!10}
Swin UNETR      & Trans.  & 82.4 & 18.3 & \checkmark  \\
\midrule
\rowcolor{utadgray!20}
TransUNet       & Hybrid  & 77.3 & 31.7 & \checkmark  \\
\rowcolor{utadgray!20}
nnFormer        & Hybrid  & 86.9 & 14.2 & ---         \\
\midrule
\rowcolor{foundorange!10}
MedSAM$^\dagger$ & Found. & 85.7 &  8.1 & \checkmark  \\
\bottomrule
\end{tabular}
}
\vspace{0.3em}

{\footnotesize
  $^\dagger$Zero-shot with bounding-box prompts.\quad Pre-train: external data required.\\
  \textcolor{cnnblue!80!black}{$\blacksquare$}\,CNN\quad
  \textcolor{transgreen!70!black}{$\blacksquare$}\,Transformer\quad
  \textcolor{utadgray!60!black}{$\blacksquare$}\,Hybrid\quad
  \textcolor{foundorange!80!black}{$\blacksquare$}\,Foundation
}

\vspace{0.5em}
\begin{notebox}[utadnavy]
  {\small
    \textbf{Auto-configured CNNs} (nnU-Net, MedNeXt, STU-Net):\quad
    $\mathbf{85.2 \pm 2.1}$\textbf{\% Dice --- 4.7 pts ahead, no pre-training}\\[2pt]
    \textbf{Transformers} (UNETR, Swin UNETR):\quad
    $\mathbf{80.5 \pm 1.8}$\textbf{\% Dice --- only reach 82\% with 5{,}050 CT pre-training}
  }
\end{notebox}
\end{pblock}

%--- Per-Organ ---------------------------------------------------------------
\begin{pblock}{Per-Organ Dice Scores \newline (from 32 of 52 studies reporting)}
\vspace{0.4em}
{\small
\begin{tabular}{@{}lcccc@{}}
\toprule
\textbf{Organ} & \textbf{Mean} & \textbf{Range} & \textbf{n} &
  \textbf{Challenge} \\
\midrule
Liver        & 94.2\% & 89.0--97.1\% & 16 & Large, well-defined  \\
Spleen       & 93.1\% & 84.0--97.0\% &  9 & Consistent shape     \\
Kidney       & 92.1\% & 78.3--98.4\% & 16 & Bilateral, variable  \\
Pancreas     & 78.1\% & 72.0--86.1\% &  8 & Small, deformable    \\
\bottomrule
\end{tabular}
}
\vspace{0.4em}
\begin{notebox}[utadred]
  {\small
    Vascular structures score \textbf{65--75\%} Dice versus
    \textbf{90\%+} for large parenchymal organs --- a known
    obstacle to surgical-navigation deployment.
  }
\end{notebox}
\end{pblock}

%--- Architecture Landscape --------------------------------------------------
\begin{pblock}{Architecture Landscape}
\vspace{0.5em}
\begin{center}
\resizebox{0.96\columnwidth}{!}{%
\begin{tikzpicture}[
  arch/.style={rectangle, rounded corners=5pt, text width=9.5cm,
               minimum height=4.0cm, align=center, font=\normalsize,
               inner sep=7pt},
  barr/.style={->, >=stealth, thick, gray!55, line width=1.5pt}]

  \node[arch, fill=cnnblue!12, draw=cnnblue, line width=1.8pt] (cnn) at (0,0) {
    {\color{cnnblue!70!black}\large\bfseries CNN}\\[3pt]
    \small nnU-Net, MedNeXt, STU-Net\\
    Dice: \textbf{82--87\%}, Params: 26--58\,M\\
    No pre-training, Fast inference
  };

  \node[arch, fill=transgreen!12, draw=transgreen, line width=1.8pt] (tr) at (11,0) {
    {\color{transgreen!60!black}\large\bfseries Transformer}\\[3pt]
    \small UNETR, Swin UNETR\\
    Dice: \textbf{79--82\%}, Params: 62--93\,M\\
    Pre-training required (5{,}050 CTs)
  };

  \node[arch, fill=foundorange!12, draw=foundorange, line width=1.8pt] (fo) at (22,0) {
    {\color{foundorange!80!black}\large\bfseries Foundation}\\[3pt]
    \small MedSAM, SAM-Med3D\\
    Dice: \textbf{76--86\%}, Params: 94\,M\\
    Interactive prompts, zero-shot capable
  };

  \draw[barr] (cnn.east) --
    node[midway, font=\footnotesize\bfseries, text=gray!60, rotate=90, fill=white, inner sep=2pt] {global context}
    (tr.west);
  \draw[barr] (tr.east) --
    node[midway, font=\footnotesize\bfseries, text=gray!60, rotate=90, fill=white, inner sep=2pt] {generalisation}
    (fo.west);
\end{tikzpicture}%
}
\end{center}

\vspace{0.4em}
{\small
\begin{tabular}{@{}lcc@{}}
\toprule
\textbf{Family} & \textbf{BTCV Dice} & \textbf{Share of 52 Studies} \\
\midrule
CNN (all)              & 74--87\% & 98.1\% \\
\quad 3D U-Net variants & 74--87\% & 82.7\% \\
Transformer (pure)     & 79--82\% & 9.6\%  \\
Hybrid (CNN+Trans.)    & 77--87\% & 7.7\%  \\
Foundation model       & 76--86\% & 3.8\%  \\
\bottomrule
\end{tabular}
}
\vspace{0.2em}
\begin{notebox}[utadnavy]
  {\footnotesize
    Architecture type matters less than training protocol: nnFormer
    (Hybrid, 86.9\%) matches STU-Net-L (CNN, 87.1\%), and
    auto-configured nnU-Net (31\,M params) rivals Swin UNETR
    (62\,M, pre-trained).
  }
\end{notebox}
\end{pblock}

%--- Survey Scope & Limitations ---------------------------------------------
\begin{pblock}{Survey Scope \& Limitations}
\small
\begin{itemize}\setlength\itemsep{1pt}
  \item \textbf{English-only search:} may under-represent
        non-English-language publications (e.g.\ Chinese-language journals)
  \item \textbf{PDF availability:} only 39.1\% of 161 AI-screened
        candidates had retrievable full-text; 98 remain paywalled
  \item \textbf{Benchmark gap:} reported Dice likely overestimates
        clinical accuracy by \textbf{5--15\%} due to controlled splits,
        test-time augmentation, and ensembling
  \item \textbf{Single-centre dominance:} only 4/52 studies (7.7\%)
        include multi-centre or prospective validation
\end{itemize}
\end{pblock}

%--- Glossary / Abbreviations -----------------------------------------------
\begin{tcolorbox}[
  title={Glossary \& Abbreviations},
  colbacktitle=utadred,
  coltitle=white,
  colback=utadgray!22,
  colframe=utadnavy!85!black,
  fonttitle=\large\bfseries,
  boxrule=1.5pt,
  arc=5pt,
  lefttitle=7pt,
  toptitle=2pt,
  bottomtitle=2pt,
  top=4pt, bottom=4pt, left=6pt, right=6pt,
  before={\vskip2pt},
  after={\vskip2pt},
  enhanced,
  drop shadow={utadnavy!25!white,opacity=0.5},
]
\footnotesize
\setlength{\tabcolsep}{3pt}
\renewcommand{\arraystretch}{1.0}
\begin{tabular}{@{}p{0.22\linewidth}p{0.72\linewidth}@{}}
\textbf{BTCV}       & Beyond Cranial Vault, 30-scan 13-organ benchmark \\
\textbf{CNN}        & Convolutional Neural Network \\
\textbf{CT / 3D}    & Computed tomography volume \\
\textbf{Dice}       & $2|A\cap B|/(|A|+|B|)$; 1.0\,=\,perfect \\
\textbf{EAU}        & European Assoc.\ of Urology \\
\textbf{HD95}       & 95th-\%ile Hausdorff dist.\ (mm) \\
\textbf{$\kappa$}   & Inter-rater agreement (Cohen's) \\
\textbf{MC-Drop.}   & Monte-Carlo Dropout \\
\textbf{nnU-Net}    & Self-config.\ U-Net (Isensee et al.) \\
\textbf{Params (M)} & Model parameters (millions) \\
\textbf{Pre-train}  & Training on external data \\
\textbf{PRISMA}     & Systematic review standard \\
\textbf{TotalSeg.}  & TotalSegmentator, 117-class nnU-Net \\
\textbf{TTA}        & Test-Time Augmentation \\
\end{tabular}
\end{tcolorbox}

\end{column}
\hfill

% ===
%  COLUMN 3 — Barriers, Recommendations, Conclusions, References
% =============================================================================
\begin{column}{0.315\textwidth}

%--- Clinical Deployment Barriers (RQ4) ------------------------------------
\begin{pblock}{Clinical Deployment Barriers (RQ4)}
\begin{enumerate}\setlength\itemsep{2pt}
  \item \textbf{Domain shift:} scanner heterogeneity degrades Dice by
        \textbf{5--15\%} across institutions; only 4/52 studies (7.7\%)
        include multi-centre validation
  \item \textbf{Vascular deficit:} 65--75\% Dice for vessels vs.\
        90\%+ for parenchymal organs
  \item \textbf{Integration complexity:} deployment requires preprocessing
        robustness, uncertainty quantification, mesh generation, and DICOM
        integration, none of which are evaluated in benchmarks
  \item \textbf{Benchmark overestimation:} published Dice likely
        overestimates clinical accuracy by \textbf{5--15\%} (controlled splits,
        TTA, ensembling)
  \item \textbf{Annotation bottleneck:} dense 3D labelling costs
        15--45 min/organ; weak/semi-supervised methods needed for scaling
\end{enumerate}
\end{pblock}

%--- Recommendations ---------------------------------------------------------
\begin{pblock}{Phased Adoption Strategy for Surgical Planning Platforms}

\begin{tcolorbox}[colback=utadnavy!10,colframe=utadnavy,arc=4pt,boxrule=1.2pt,
  top=4pt,bottom=4pt,left=7pt,right=7pt,before={\vskip0pt},after={\vskip0pt}]
  {\normalsize\bfseries\color{utadnavy!80!black} 1. Baseline Segmentation}\\[2pt]
  {\small TotalSegmentator or nnU-Net for robust multi-organ segmentation
  without site-specific tuning}
\end{tcolorbox}
\begin{tcolorbox}[colback=utadnavy!10,colframe=utadnavy,arc=4pt,boxrule=1.2pt,
  top=4pt,bottom=4pt,left=7pt,right=7pt,before={\vskip0pt},after={\vskip0pt}]
  {\normalsize\bfseries\color{utadnavy!80!black} 2. Uncertainty-Guided QA}\\[2pt]
  {\small MC-Dropout / TTA disagreement, Flag low-confidence cases for
  expert review}
\end{tcolorbox}
\begin{tcolorbox}[colback=utadnavy!10,colframe=utadnavy,arc=4pt,boxrule=1.2pt,
  top=4pt,bottom=4pt,left=7pt,right=7pt,before={\vskip0pt},after={\vskip0pt}]
  {\normalsize\bfseries\color{utadnavy!80!black} 3. Interactive Refinement}\\[2pt]
  {\small MedSAM with bounding-box prompts, Surgeon-guided correction of
  uncertain regions}
\end{tcolorbox}

\vspace{0.4em}
\begin{notebox}[utadnavy]
  {\small
    \emph{Implementation quality matters more than
    architectural novelty.} nnU-Net without modifications outperforms many
    specialised methods on standard benchmarks.
  }
\end{notebox}
\end{pblock}

%--- Conclusions -------------------------------------------------------------
\begin{pblock}{Conclusions}
\begin{enumerate}\setlength\itemsep{2pt}
  \item \textbf{CNNs dominate} (98.1\%); nnU-Net is the strongest reported
        baseline (82.3\% Dice on BTCV, no pre-training)
  \item \textbf{Auto-CNNs outperform transformers} on BTCV:
        $85.2 \pm 2.1$\% vs.\ $80.5 \pm 1.8$\%;
        MedNeXt (86.2\%) leads without any pre-training
  \item \textbf{Transformers are complementary:} Swin UNETR (82.4\%) matches
        nnU-Net when pre-trained on 5,050 CT volumes
  \item \textbf{Foundation models} enable interactive refinement
        (MedSAM 85.7\% fine-tuned); zero-shot ($\sim$76\%) still trails
        supervised methods
  \item \textbf{Phased adoption} (TotalSegmentator $\to$ QA $\to$ MedSAM)
        is the recommended strategy for surgical planning platforms
\end{enumerate}
\end{pblock}

%--- Key References ----------------------------------------------------------
\begin{pblock}{References}
{\footnotesize
\begin{enumerate}\setlength\itemsep{2pt}
  \item[{[1]}] Isensee et al.\ \emph{nnU-Net.} Nat.\ Methods, 2021.
  \item[{[2]}] Tang et al.\ \emph{Swin UNETR.} MICCAI WS, 2022.
  \item[{[3]}] Wasserthal et al.\ \emph{TotalSegmentator.} Radiol.\ AI, 2023.
  \item[{[4]}] Ronneberger et al.\ \emph{U-Net.} MICCAI, 2015.
  \item[{[5]}] \c{C}i\c{c}ek et al.\ \emph{3D U-Net.} MICCAI, 2016.
  \item[{[6]}] Ma et al.\ \emph{Segment Anything in Medical Imaging.} Nat.\ Commun., 2024.
  \item[{[7]}] Roy et al.\ \emph{MedNeXt.} MICCAI, 2023.
  \item[{[8]}] Hatamizadeh et al.\ \emph{UNETR.} WACV, 2022.
\end{enumerate}
}
\end{pblock}

\end{column}
\end{columns}
\end{frame}
\end{document}
